\documentclass[a4paper,10pt]{report}\input{../0.Préambule/Préambule_cours_prof}\begin{document}

\newpage
\pagestyle{empty}
\Large{Troisième Générale \hfill Session Avril 2022

\vspace{1cm}
\begin{minipage}{8cm}
Nom : \dotfill
\end{minipage}
\hspace{5mm}
\begin{minipage}{8cm}
Prénom : \dotfill
\end{minipage}

\vspace{0.5cm}
\noindent\hrulefill \vspace{10mm}
\begin{center}
\textbf{\Huge{Brevet Blanc de Mathématiques}}
\end{center}\vspace{10mm}
\noindent\hrulefill


\vspace{0.2cm}

\begin{center}
\subsection*{\hspace{4.5cm} Durée de l'épreuve : 2 heures\hspace{4.5cm} }
\end{center}

\vspace{0.2cm}


\noindent \textbf{Le sujet est composé de $8$ pages.}

\medskip
\begin{enumerate}
\item[•] Toute réponse doit être parfaitement justifiée (sauf indication contraire dans la consigne).
\item[•] La qualité, la clarté de la rédaction et la maîtrise de la langue seront prises en compte dans l'appréciation de la copie.
\item[•] Les exercices seront traités dans l'ordre de votre choix.
\item[•] L'emploi de la calculatrice est autorisé.
\end{enumerate}

\vspace{0.2cm}

\noindent \textbf{Barème des exercices :}

\begin{center}
\begin{minipage}{15cm}
\renewcommand{\arraystretch}{1.2}
\begin{tabularx}{\linewidth}{|*{4}{>{\centering \arraybackslash}X|}}
\hline
\rowcolor{gray!20}Exercice & \multicolumn{2}{|c|}{Nombre de points} & Compétences \\\hline
$1$ & & $14$ & Calculer \\\hline
$2$ & & $12$ & Calculer \\\hline
$3$ & & $15$ &  \\\hline
$4$ & & $16$ &  \\\hline
$5$ & & $14$ & Raisonner \\\hline
$6$ & & $9$ & Chercher \\\hline
$7$ & & $12$ &  \\\hline
$8$ & & $8$ & Algorithme \\\hline
Total & & $100$ & \multicolumn{1}{c}{}\\\cline{1-3}
\end{tabularx}
\end{minipage}
\end{center}

\vspace{0.5cm}

\begin{center}
Collège Charles DULLIN
\end{center}

\newpage
\textbf{Exercice 1}: \hfill \textbf{14 points}
\medskip

Pour les questions $A$, $B$ et $C$, on écrira \textbf{toutes les étapes de calcul} menant à la réponse sur sa copie. Il ne sera attribué aucun points si la réponse n'est pas justifiée soigneusement

\renewcommand{\arraystretch}{2.2}
\begin{tabularx}{\linewidth}{|m{4mm}|m{6cm}|*{3}{>{\centering \arraybackslash}X|}}
\cline{3-5}
\multicolumn{2}{c|}{} & Réponse n \degre 1 & Réponse n \degre 2 & Réponse n \degre 3  \\\hline
$A$ & Le nombre $\dfrac{7}{5} - \dfrac{4}{5} \times \dfrac{1}{3}$ & $\dfrac{17}{15}$ & $\dfrac{1}{5}$ & $-\dfrac{3}{10}$ \\\hline
$B$ & Le nombre $\dfrac{4^{-8} \times 4^{5}}{4^3}$ & $4^{-6}$ & $1$ & $4^{-16}$ \\\hline
$C$ & L'équation $5x+12 = 3$ a pour solution & $1,8$ & $-3$ & $-1,8$ \\\hline
$D$ & La formule à saisir dans la cellule $B2$ avant de la recopier vers le bas est : 

\begin{center}
\begin{tableur}{3}
\hline
&\cellcolor{white}{Valeurs} & \cellcolor{white}{$3x^2-8x$}& \cellcolor{white}{}\\\hline
&\cellcolor{white}{$0$} & \cellcolor{white}{$0$}& \cellcolor{white}{}\\\hline
&\cellcolor{white}{$1$} & \cellcolor{white}{$-5$}& \cellcolor{white}{}\\\hline
&\cellcolor{white}{$2$} & \cellcolor{white}{$-4$}& \cellcolor{white}{}\\\hline
&\cellcolor{white}{$3$} & \cellcolor{white}{$3$}& \cellcolor{white}{}\\\hline
&\cellcolor{white}{$4$} & \cellcolor{white}{$16$}& \cellcolor{white}{}\\\hline
&\cellcolor{white}{$5$} & \cellcolor{white}{$35$}& \cellcolor{white}{}\\\hline
\end{tableur}
\end{center}
\arrayrulecolor{black}
& $3 ^{*}A2 \hat{\quad}2-8 ^{*}A2$ & $=3 ^{*}A1 \hat{\quad}2-8 ^{*}A1$ & $=3 ^{*}A2 \hat{\quad}2-8 ^{*}A2$  \\\hline
\end{tabularx}

\medskip
\textbf{Exercice 2}: \hfill \textbf{12 points}
\medskip

On considère l'expression suivante : $B=(5x-1)(2x-3) + (5x-1)(3x-4)$
\begin{enumerate}
\item Développer, réduire et ordonner l'expression $B$.
\item Factoriser l'expression $B$
\item Parmi les trois expression précédente, choisir l'expression qui vous parait la plus adaptée et calculer la valeur de $B$ pour $x=\dfrac{1}{5}$.
\end{enumerate}

\newpage
\textbf{Exercice 3}: \hfill \textbf{15 points}
\medskip

Un professeur de SVT demande aux $29$ élèves d'une classe de sixième de faire germer des graines de blé chez eux. Le tableau Excel ci-dessous donne les tailles des plantules (petites plantes) des $29$ élèves à $10$ jours après la mise en germination.

\begin{center}
\renewcommand{\arraystretch}{1.5}
\begin{tableur}{13}
\hline
& Tailles des plantules (en cm) & $0$ & $8$ & $12$ & $14$ & $16$ & $17$ & $18$ & $19$ & $20$ & $21$ & $22$ & Total\\\hline
& Effectifs & $1$ & $2$ & $2$ & $4$ & $2$ & $2$ & $3$ & $3$ & $4$ & $4$ & $2$ & \\\hline
\end{tableur}
\end{center}

\begin{enumerate}
\item Combien de plantules ont une taille qui mesure au plus $12$cm ?
\item Donner l'étendue de cette série.
\item Calculer la moyenne de cette série. \textbf{Arrondir au dixième près}.
\item Déterminer la médiane de cette série et interpréter le résultat.
\item On considère qu'un élève a bien respecté le protocole si la taille de la plantule à $10$ jours est supérieure ou égale à $14$cm. 

Quel pourcentage des élèves de la classe a bien respecté le protocole ? \textbf{Arrondir à l'unité}.
\item Le professeur a fait lui-même la même expérience en suivant le même protocole. Il a relevé la taille obtenue à $10$ jours de germination.

Prouver que, si on ajoute la donnée du professeur à cette série, la médiane ne changera pas.
\end{enumerate}


\textbf{Exercice 4}: \hfill \textbf{16 points}
\medskip

\begin{enumerate}
\item Décomposer les nombres $162$ et $108$ en produits de facteurs premiers.
\item Déterminer deux diviseurs communs  aux nombres $162$ et $108$ plus grands que $10$.
\item Un snack vend des barquettes composées de nems et de samossas. Le cuisinier a préparé $162$ nems et $108$ samossas. Dans chaque barquette le nombre de nems et de samossas doit être le même.

Tous les nems et tous les samossas doivent être utilisés.
\begin{enumerate}
\item Le cuisiner peut-il réaliser $36$ barquettes ?
\item Quel nombre maximal de barquettes pourra-t-il réaliser ?
\item Dans ce cas, combien y aura-t-il de nems et de samossas dans chaque barquette ?
\end{enumerate}
\end{enumerate}

\textbf{Exercice 5}: \hfill \textbf{14 points}
\medskip

Aya participe à un rallye VTT sur un parcours balisé. \\Le trajet est représenté en traits pleins sur la figure ci-dessous. \\Le départ du rallye est en $A$ et l'arrivée est en $G$.

\noindent \begin{minipage}{8cm}
\begin{center}
\begin{tikzpicture}[line cap=round,line join=round,>=triangle 45,x=1.0cm,y=1.0cm,scale=0.8]
\clip(0.2,-0.8) rectangle (10,6.8);
\draw [line width=1.2pt] (2.,6.)-- (5.,6.);
\draw [line width=1.2pt,dash pattern=on 3pt off 3pt] (5.,6.)-- (6.,6.);
\draw [line width=1.2pt,dash pattern=on 3pt off 3pt] (6.,0.)-- (8.,0.);
\draw [line width=1.2pt] (8.,0.)-- (9.5,0.);
\draw [line width=1.2pt] (8.,0.)-- (5.,6.);
\draw [line width=1.2pt,dash pattern=on 3pt off 3pt] (6.,6.)-- (6.,0.);
\begin{scriptsize}
\draw [fill=black] (2,6.) circle (1.0pt);
\draw[color=black] (1.64,6.35) node {\large{$A$}};
\draw [fill=black] (5.,6.) circle (1.0pt);
\draw[color=black] (4.64,6.37) node {\large{$B$}};
\draw[color=black] (2.84,6.47) node {\large{$7$km}};
\draw [fill=black] (6.,6.) circle (1.0pt);
\draw[color=black] (6.28,6.09) node {\large{$C$}};
\draw[color=black] (5.54,6.47) node {\large{$1,5$km}};
\draw [fill=black] (6.,0.) circle (1.0pt);
\draw[color=black] (5.66,-0.39) node {\large{$E$}};
\draw [fill=black] (8.,0.) circle (1.0pt);
\draw[color=black] (8.02,-0.39) node {\large{$F$}};
\draw [fill=black] (9.5,0.) circle (1.0pt);
\draw[color=black] (9.5,-0.43) node {\large{$G$}};
\draw[color=black] (8.7,0.43) node {\large{$3,5$km}};
\draw[color=black] (6.55,4.85) node {\large{$2$km}};
\draw[color=black] (5.3,1.51) node {\large{$5$km}};
\draw [fill=black] (6.,4.) circle (1.0pt);
\draw[color=black] (5.5,3.81) node {\large{$D$}};
\end{scriptsize}
\end{tikzpicture}
\end{center}
\end{minipage}
\hspace{10mm}
\begin{minipage}{8cm}
\textit{Les points $A$, $B$, $C$ sont alignés}\\
\textit{Les points $C$, $D$, $E$ sont alignés}\\
\textit{Les points $E$, $F$, $G$ sont alignés}\\
\textit{Les points $B$, $D$, $F$ sont alignés}\\

\medskip
\textit{$BCD$ est rectangle en $C$.}\\
\textit{$DEF$ est rectangle en $E$.}
\end{minipage}

\begin{enumerate}
\item Montrez que $BD = 2,5$km
\item Justifiez que les droites $(BC)$ et $(EF)$ sont parallèles.
\item En déduire la longueur $DF$.
\item Calculez la longueur totale du parcours effectué par Aya lors de ce rallye.
\end{enumerate}

\textbf{Exercice 6}: \hfill \textbf{9 points}
\medskip

C'est bientôt le printemps, Mme Durand doit changer les quatre pneus de sa voiture. Elle souhaite commander ses nouveaux pneus sur un site internet qui propose une réduction de $20 \%$ sur chaque pneu acheté de la marque Goodyear.

\noindent A l'aide des informations suivantes, déterminer la somme qu'elle devra payer pour l'achat de ses quatre pneus si elle profite de l'offre de réduction.

\begin{doc}[Résultat d'un test]

Mme Durand, à l'aide des informations que vous nous avez fournies concernant votre véhicule, voici les dimensions des pneus qui conviennent pour votre voiture.

\begin{center}
\renewcommand{\arraystretch}{1}
\arrayrulecolor{black}
\begin{tabularx}{5cm}{|>{\columncolor{lightgray!50}\centering \arraybackslash}X|>{\centering \arraybackslash}X|}
\hline
Diamètre & $15$ \\\hline
Hauteur & $65$ \\\hline
Charge & $88$ \\\hline
Largeur & $185$ \\\hline
Vitesse & $H$ \\\hline
\end{tabularx}
\end{center}
\end{doc}

\begin{doc}[Ce que signifie les information sur un pneu :]
\begin{center}
\includegraphics[scale=0.6]{../40.Image/DC2 2.png}
\end{center}
\end{doc}

\begin{doc}[Les tarifs affichés sur le site internet]
\begin{center}
\includegraphics[scale=0.6]{../40.Image/DC2 1.png}
\end{center}
\end{doc}

\newpage
\textbf{Exercice 7}: \hfill \textbf{12 points}
\medskip

\medskip

\noindent \begin{minipage}{7cm}
Pour la course à pied en montagne, certains sportifs mesurent leur
performance par la \textbf{vitesse ascensionnelle}, notée $V_a$.

\vspace{5mm}
$V_a$ est le quotient du dénivelé de la course, exprimé en mètres, par la durée, exprimée en heure. \hfill
\end{minipage}
\hspace{5mm}
\begin{minipage}{9cm}
\begin{center}
\includegraphics[scale=0.5]{../40.Image/DC2 3.png}
\end{center}
\end{minipage}

\medskip

\noindent Par exemple : pour un dénivelé de \np{4500} m et une durée de parcours de 3 h : $V_a = \np{1500}$ m/h.

\medskip
\noindent \textbf{Rappel:} le dénivelé de la course est la différence entre l'altitude à l'arrivée et l'altitude au départ.

\medskip

Un coureur de haut niveau souhaite atteindre une vitesse ascensionnelle d'au moins \np{1400} m/h lors de sa prochaine course.

\begin{center}
\emph{La figure ci-dessous n'est pas représentée en vraie grandeur.}
\includegraphics[scale=0.5]{../40.Image/DC2 4.png}
\end{center}

Le parcours se décompose en deux étapes (voir figure 2) :

\medskip

\setlength\parindent{12mm}
\begin{itemize}
\item[$\bullet~~$]Première étape de \np{3800} m pour un déplacement horizontal de \np{3790} m.
\item[$\bullet~~$]Seconde étape de $4,1$ km avec un angle de pente d'environ $12$\degres.
\end{itemize}
\setlength\parindent{0mm}

\bigskip

\begin{enumerate}
\item Vérifier que le dénivelé de la première étape est environ $275,5$ m.
\item Quel est le dénivelé de la seconde étape ?
\item Depuis le départ, le coureur met $48$ minutes pour arriver au sommet.

Le coureur atteint-il son objectif ?
\end{enumerate}

\newpage
\textbf{Exercice 8}: \hfill \textbf{8 points}
\medskip

\noindent Gabrielle a crée quatre variables \textbf{x}, \textbf{Étape 1}, \textbf{Étape 2} et \textbf{Résultat}, puis elle a saisi le script ci-dessous (en couleur sur l'\textbf{ANNEXE}) en langage Scratch.
\medskip

\begin{center}
\begin{scratch}[print,fill blocks]
\blockinit{quand \greenflag est cliqué}
\blockmove{demander \ovalnum{Valeur de x ?} et attendre}
\blockvariable{mettre {\selectmenu{x} à \ovalsensing{réponse}}}
\blockvariable{mettre {\selectmenu{Étape 1} à \ovaloperator{\ovalvariable{x} + \ovalnum{4}}}}
\blockvariable{mettre {\selectmenu{Étape 2} à \ovaloperator{\ovaloperator{\ovalnum{2} *  \ovalvariable{x}} - \ovalnum{3}}}}
\blockvariable{mettre {\selectmenu{Résultat} à \ovaloperator{ \ovalvariable{Etape 1} * \ovalvariable{Etape 2}}}}
\blocklook{dire \ovaloperator{regrouper \ovalnum{Le résultat est } et \ovalvariable{Résultat}}}
\end{scratch}
\end{center}

\begin{enumerate}
\item Vérifier que si la valeur saisie de \textbf{x} est $5$, alors le résultat est $63$
\item Quel résultat obtient-on si la valeur de \textbf{x} est $-3$ ?
\item Parmi les expression suivantes, recopier celle que correspond au résultat affiché par le script.
\begin{enumerate}
\item[•] $A= (x+4)\times (2x-3)$
\item[•] $B= x+4\times 2x-3$
\item[•] $C= x+4\times (2x-3)$
\end{enumerate}
\item Pour quelle(s) valeur(s) de $x$ obtient-on un résultat égal à $0$ ?
\end{enumerate}

}
\end{document}
